% !TEX program = xelatex
% AI-effects 프로젝트 — reference/ 소장 문헌 중 AI 영향을 "노출도(exposure)"가 아니라
% (1) 실제 채택(adoption) 여부, (2) 실제 사용(usage) 양·패턴으로 정의한 실증연구를
% 데이터 출처와 방법론을 중심으로 정리한 작업 문헌(report/).
% 01_AI영향_실증연구_정리.tex의 §4(채택)·§5(실사용)를 씨줄로 삼아 노출도 절을 제외하고
% 데이터 출처의 기관적 성격과 식별전략(방법론)을 심화한다.
% 컴파일: xelatex -> biber -> xelatex -> xelatex (kotex, biblatex/biber 필요)
\documentclass[11pt]{article}

\usepackage{fontspec}
\usepackage{kotex}
\usepackage[a4paper,margin=2.2cm]{geometry}
\usepackage{longtable}
\usepackage{booktabs}
\usepackage{array}
\usepackage{amsmath}
\usepackage{xcolor}
\usepackage{hyperref}
\usepackage{enumitem}
\usepackage{titlesec}
\usepackage[backend=biber,style=authoryear,maxcitenames=2,uniquelist=false,sorting=nyt]{biblatex}
\addbibresource{../../reference/references.bib}

\hypersetup{colorlinks=true,linkcolor=blue!50!black,citecolor=blue!50!black,urlcolor=blue!50!black}
\newcolumntype{L}[1]{>{\raggedright\arraybackslash}p{#1}}
\setlength{\LTleft}{0pt}
\setlength{\LTright}{0pt}
\renewcommand{\arraystretch}{1.15}
\setlength{\tabcolsep}{4pt}

\title{AI 채택(Adoption)과 실사용(Usage) 기반 실증연구:\\[1mm]
데이터 출처와 방법론\\[2mm]
\large --- 노출도(exposure) 지표를 제외한, 실현된 AI 영향의 측정 ---}
\author{AI-effects 프로젝트}
\date{\today}

\begin{document}
\maketitle

\begin{abstract}
\noindent 본 문헌은 AI(인공지능)의 거시경제·노동시장 영향을 사전적 노출(ex-ante exposure) 지표가 아니라 (1) 기업·근로자의 \textbf{실제 AI 채택(adoption)} 여부·정도, (2) AI가 \textbf{실제로 얼마나·어떻게 사용되는지(usage)}로 정의한 실증연구 28편을, 각 연구가 활용한 \emph{데이터 출처}(수집기관, 조사설계·표본, 시점·주기, 원자료의 성격)와 \emph{방법론}(핵심변수 정의, 추정모형, 식별전략)을 중심으로 정리한다. 노출도 기반 연구는 이미 \texttt{01\_AI영향\_실증연구\_정리.tex}·\texttt{02\_AI노출도\_지표\_비교.tex}·\texttt{04\_한국AI실증연구\_방법론계보.tex}·\texttt{05\_Felten계승연구\_직업분류매칭여부.tex}에서 다루었으므로 본 문헌에서는 제외한다. 신규로 \citet{cheon-2025-ai-exposure-vs-ai-adoption}의 AI Firm Exposure(AIFE, AI 공급기업의 서비스 텍스트를 직업 과업과 LLM으로 매칭해 산출한 "실제 적용도" 지표)를 채택 범주에 편입해, 기존 report에는 반영되지 않았던 이 핵심 문헌을 처음으로 상세히 다룬다.
\end{abstract}

\tableofcontents

% ============================================================
\section{서론}
% ============================================================

\texttt{CLAUDE.md}가 규정하는 이 저장소의 목표는 AI exposure(사전적·이론적 노출)와 AI usage(실현된 실제 사용) 데이터를 결합해, 노출만으로는 포착하지 못하는 실제 채택·활용 효과를 분석하는 것이다. \texttt{01\_AI영향\_실증연구\_정리.tex}는 이 목표에 필요한 선행연구 지형을 노출도·채택·실사용 세 갈래로 개관했고, 각 연구를 (1)정의·지표 (2)데이터 (3)종속변수 (4)국가·범위 네 칸의 압축 표로 요약했다. 그러나 본 프로젝트가 \texttt{data/proc/merged/}를 실제로 설계하려면, 요약표의 "데이터" 한 칸에 압축된 정보—누가 그 자료를 수집했는지, 표본이 반복측정되는지 단발성인지, 자기보고인지 관측된 로그인지—와 "방법론"이 실제로 얼마나 강한 인과추론을 뒷받침하는지가 훨씬 중요해진다. 본 문헌은 노출도 절을 제외하고 \textbf{채택(adoption)}과 \textbf{실사용(usage)} 두 갈래만을 대상으로, 데이터 출처와 방법론을 씨줄로 삼아 다시 정리한다.

이 재정리가 특히 필요한 이유는 \citet{cheon-2025-ai-exposure-vs-ai-adoption}(전병유·신영민, 2025)이 제시한 다음 사실 때문이다: 한국형 AIOE(기술적 노출)와 AI Firm Exposure(AIFE, 실제 시장에서 팔리고 있는 AI 서비스가 어떤 직업의 과업을 대상으로 하는지 LLM으로 측정한 지표) 간 피어슨 상관계수는 약 0.35에 불과하다. 이는 "기술적으로 가능한 일"과 "실제로 시장에서 실현되고 있는 일"이 별개의 정보를 담고 있음을 뜻하며, 노출도 지표만으로 구성된 분석은 이 괴리를 볼 수 없다. 이 문헌은 \texttt{log.md}(2026-09-16)에 "아직 미반영"으로 남아 있던 이 논문을 채택(adoption) 범주의 핵심 사례로 처음 편입한다—AIFE는 근로자·기업의 자기보고(수요측)가 아니라 AI 공급기업의 실제 제품·서비스 텍스트(공급측)에서 직접 도출된다는 점에서, 아래 \S\ref{sec:adoption-data}에서 다룰 다른 모든 채택 지표와 성격이 다르다.

정리 대상은 총 28편이다: 채택 기반 16편(\S\ref{sec:adoption}), 실사용 기반 12편(\S\ref{sec:usage}). 일부 논문(\citealp{nam-2026-ai-macroeconomic-impact-analysis, kiet-2024-ai-labor-market-industry, chang-2024-ai-development-employment-effects, noh-2025-ai-labor-coexistence, noh-2025-ai-based-manufacturing-innovation-employment})은 노출도와 채택을 한 보고서 안에서 병행하므로, 본 문헌에서는 채택과 직접 관련된 장·절만 추출해 다룬다(노출도 부분은 \texttt{04\_한국AI실증연구\_방법론계보.tex} 참조).

% ============================================================
\section{정리 기준}
% ============================================================

각 연구는 다음 두 축으로 정리한다.

\begin{enumerate}[label=(\arabic*)]
  \item \textbf{데이터 출처}: 누가/어느 기관이 수집했는가(정부통계·자체 서베이·행정DB·텍스트마이닝·AI기업 대화로그 등), 조사설계·표본추출 방법, 표본규모, 조사시점 및 주기(1회성 단면 vs. 반복 단면 vs. 패널), 원자료의 성격(자기보고 vs. 관측 행정자료 vs. 텍스트/로그).
  \item \textbf{방법론}: 처치변수·핵심변수의 정의, 추정모형과 식별전략(단순 기술통계 $\to$ 회귀/고정효과 $\to$ 준실험(DID·사건연구·IV) $\to$ 구조모형/실험), 강건성 검증 방식.
\end{enumerate}

\S\ref{sec:adoption}은 채택 기반 연구를, \S\ref{sec:usage}는 실사용 데이터 기반 연구를 다루며, 각 절은 데이터 출처 유형별 개관, 방법론 유형별 개관, 연구별 요약표(3열: 문헌·데이터 출처·방법론) 순으로 구성한다. \S\ref{sec:compare}는 두 갈래를 데이터 출처·방법론 측면에서 직접 비교하고, \S\ref{sec:conclusion}은 본 프로젝트에 대한 시사점을 논의한다.

% ============================================================
\section{채택(Adoption) 기반 실증연구}
\label{sec:adoption}
% ============================================================

\subsection{데이터 출처 유형}
\label{sec:adoption-data}

채택 기반 연구의 데이터 출처는 크게 일곱 갈래로 나뉜다.

\begin{itemize}[leftmargin=*]
  \item \textbf{정부 기업패널조사(반복측정)}: 통계청(현 국가데이터처) 기업활동조사와 한국노동연구원 사업체패널조사(WPS)가 핵심 자료원이다. 기업활동조사는 매년, WPS는 격년으로 실시되는 반복 패널이며, AI 도입 여부는 기업의 자기보고 더미로 조사된다. \citet{kiet-2021-ai-adoption-productivity}(2017--2019)부터 \citet{kim-2025-ai-adoption-firm-performance}(2017--2023), \citet{nam-2026-ai-macroeconomic-impact-analysis}(2017--2023)까지 동일 조사가 시점을 넓혀가며 반복 활용되고, \citet{lee-2025-ai-adoption-determinants-performance}는 2023년 WPS 10차에 OECD(2023) 국제표준 문항을 최초로 추가했다는 점에서 국제비교 가능성을 획득한 시점이다.
  \item \textbf{단면 실태조사(전수 또는 표본 서베이)}: \citet{kim-2021-ai-adoption-status-major-industries}(KISDI, 한국갤럽 면대면조사)와 \citet{spri-2024-ai-industry-survey}(과기정통부·SPRI, 국가승인통계 전수조사)는 반복되지 않는 단면 조사로, 인과모형 없이 도입률 자체를 기술통계로 보고한다.
  \item \textbf{텍스트마이닝 기반 간접식별}: \citet{seo-2025-ai-adoption-characteristics-business-reports}는 금융감독원 DART 사업보고서 전수(2010--2024)에서 AI 키워드 등장 여부로 도입을 간접 식별한다. 이는 설문 응답이 아니라 기업이 이미 법적으로 공시한 텍스트에서 도입 여부를 추출한다는 점에서, 응답 편향 없이 장기 시계열(15년)을 확보할 수 있는 독특한 자료원이다.
  \item \textbf{구인공고 기반 간접식별(Babina et al.\ 2020 방법론)}: \citet{chang-2024-ai-development-employment-effects}(4장)와 \citet{noh-2025-ai-based-manufacturing-innovation-employment}는 사람인·잡코리아 온라인 구인공고에서 AI 관련 숙련 수요가 나타나는 사업체를 "AI 도입 기업"으로 간주하고, 이를 고용보험DB와 결합한다. 도입 여부를 기업의 자기보고가 아니라 채용 행위(구인공고)라는 관측 가능한 행동으로 대리(proxy)한다는 점이 특징이다.
  \item \textbf{국제 근로자·기업 서베이}: \citet{bick-2026-mind-the-gap-ai-adoption}는 미국·유럽 6개국 근로자 55,415명 대상 자체 서베이(Bilendi 패널)와 Eurostat EU-ICT-Firm(32개국 157,000개 기업), 미국 Census BTOS를 결합해 국가 간 도입률을 비교한다.
  \item \textbf{서베이 기반 과업단위 채택지수}: \citet{bick-2026-what-work-does-generative}는 Real-Time Population Survey(RPS)로 O*NET DWA(세부업무활동) 수준의 genAI 채택률을 직접 서베이로 측정한 신규 모듈을 구축했다—기업·근로자 이분법적 도입 여부가 아니라 \emph{과업 단위} 채택 강도를 서베이로 포착한 최초 사례다.
  \item \textbf{AI 공급기업 서비스 텍스트 LLM 매칭(공급측 실현 지표)}: \citet{cheon-2025-ai-exposure-vs-ai-adoption}의 AIFE는 위 여섯 갈래와 근본적으로 다르다. 수요측(기업·근로자)의 자기보고나 행동을 관측하는 것이 아니라, 공급측인 AI 기업 355개(대표기업 51개+스타트업 339개)가 실제로 판매하는 서비스의 텍스트 내용을 GPT로 직업별 과업 텍스트와 매칭해 "이 직업의 과업을 대상으로 하는 AI 서비스가 실제로 시장에 존재하는가"를 직접 측정한다. Fenoaltea et al.(2024)의 AI Startup Exposure(AISE) 방법론을 한국에 처음 적용한 사례다.
\end{itemize}

로봇 도입(비AI, 준거 사례)의 \citet{humlum-2019-robot-adoption-and-labor-market}는 관세기록(로봇 직접수입)과 통계청격 행정DB(덴마크 IDA/FIDA/FirmStat)를 결합한 행정자료형 데이터로, 자기보고에 의존하지 않는 도입 이벤트 식별의 모범사례로 함께 다룬다. \citet{noh-2025-ai-labor-coexistence}(3--9장)는 KBS·한국노동연구원·경총 공동설문과 산업별 질적 인터뷰로 도입·활용 정도를 자기보고와 인터뷰로 병행 관측한다.

\subsection{방법론 유형}

채택 기반 연구의 식별전략은 자료의 반복측정 가능성에 비례해 정교해지는 경향이 뚜렷하다.

\begin{itemize}[leftmargin=*]
  \item \textbf{순수 기술통계}: \citet{kim-2021-ai-adoption-status-major-industries}, \citet{spri-2024-ai-industry-survey}, \citet{seo-2025-ai-adoption-characteristics-business-reports}는 단면 자료의 한계를 인정하고 인과모형 없이 도입률·비율만 보고한다.
  \item \textbf{프로빗/단순회귀}: \citet{kiet-2021-ai-adoption-productivity}·\citet{kiet-2024-ai-adoption-firms-policy}는 AI 도입의 결정요인을 프로빗으로, \citet{chang-2024-ai-development-employment-effects}(4장)는 도입-신규채용 관계를 OLS로 추정한다.
  \item \textbf{패널 고정효과(TWFE)}: \citet{kiet-2024-ai-adoption-firms-policy}·\citet{kim-2025-ai-adoption-firm-performance}는 이원고정효과 패널회귀로 도입의 성과효과를 추정한다.
  \item \textbf{준실험: 사건연구·다시점DID}: \citet{kim-2025-ai-adoption-firm-performance}는 사건연구(event-study)로 parallel trends를 검증하고, \citet{lee-2025-ai-adoption-determinants-performance}와 \citet{kiet-2024-ai-labor-market-industry}(3장)는 Callaway and Sant'Anna(2021) 다기간 DID로 코호트별 처치효과(ATT)를 식별한다.
  \item \textbf{도구변수(IV)·준실험 결합}: \citet{kiet-2021-ai-adoption-productivity}는 동종산업 AI활용률·4차산업혁명 인프라 더미를 도구변수로, \citet{nam-2026-ai-macroeconomic-impact-analysis}는 leave-one-out 도입률을 도구변수로 한 패널IV와 PSM-DID를 병행하며, \citet{chang-2024-ai-development-employment-effects}(5장)는 IV(2SLS)와 Entropy Balancing(Czarnitzki et al.\ 2023)을 결합한다.
  \item \textbf{구조적 동태모형}: \citet{humlum-2019-robot-adoption-and-labor-market}는 Rust(1987)식 동태적 이산선택 도입모형과 Lee-Wolpin(2006) 동태적 직종선택모형을 결합한 일반균형 구조모형으로 도입비용과 로봇세 정책효과까지 시뮬레이션한다—채택 기반 연구 중 인과추론 강도가 가장 높은 사례다.
  \item \textbf{국가 간 분해(Oaxaca-Blinder)}: \citet{bick-2026-mind-the-gap-ai-adoption}는 미국--유럽 도입률 격차를 구성요인과 계수차이로 분해한다.
  \item \textbf{LLM 텍스트 매칭}: \citet{cheon-2025-ai-exposure-vs-ai-adoption}의 AIFE와 \citet{bick-2026-what-work-does-generative}의 채택지수 구축은 회귀분석이 아니라 지표 구축 자체가 방법론적 기여이며, 구축된 지표를 이용한 상관분석·설명력 비교(후자는 2SLS 학습채널 식별까지 포함)로 보완된다.
\end{itemize}

\subsection{연구별 요약}

\begin{longtable}{L{2.5cm}L{6.4cm}L{6.4cm}}
\caption{채택(Adoption) 기반 실증연구: 데이터 출처와 방법론}\label{tab:adoption}\\
\toprule
\textbf{문헌} & \textbf{데이터 출처} & \textbf{방법론}\\
\midrule
\endfirsthead
\multicolumn{3}{l}{\small (표 \thetable\ 계속)}\\
\toprule
\textbf{문헌} & \textbf{데이터 출처} & \textbf{방법론}\\
\midrule
\endhead
\bottomrule
\endfoot

\citet{humlum-2019-robot-adoption-and-labor-market} &
덴마크 IDA/FIDA/FirmStat(1995--2016) 매칭 고용주-근로자 행정DB + 관세기록 UHDI(HS 847950, 1993--2015) + VITA 로봇도입 설문(2018, 응답률 97\%)을 결합, 5단계 필터링으로 로봇도입 이벤트 454건(293개 기업) 식별 &
직종별 차별적 생산성효과를 반영한 CES 생산함수 기반 Rust(1987)식 동태적 이산선택 도입모형; 매칭(Exact-Mahalanobis) 이벤트스터디/DID로 기술파라미터 GMM 추정($\sigma$=0.493); Lee-Wolpin(2006) 동태적 직종선택+CCP추정량과 결합한 일반균형모형, MSM으로 로봇세 정책실험\\

\citet{kiet-2021-ai-adoption-productivity} &
통계청 기업활동조사+광업제조업조사 연계 패널(2017--2019, 제조업, N=12,599) + 제조업 기업 503개 자체 실태조사(2021) &
AI도입 결정요인 프로빗(산업·연도FE); 생산성효과는 OLS/IV(도구변수: 동종산업 AI활용률, BIM 인프라 더미, 1단계 F=72.59), 종속변수 ln(부가가치/종업원)\\

\citet{kiet-2024-ai-adoption-firms-policy} &
통계청 기업활동조사 패널(2017--2022) + AI활용 실태조사(560개 기업) &
결정요인 프로빗(기업규모·무형자산·기업연령·R\&D집약도); 생산성효과는 이원고정효과(TWFE) 패널회귀, AI$\times$클라우드/3D프린팅 교호항으로 보완성 분석\\

\citet{kim-2021-ai-adoption-status-major-industries} &
5대산업 종사자 20인이상 기업 368개, 한국갤럽 면대면(일부 온라인) 할당추출 실태조사(2021.4--5) &
순수 기술통계(도입률·산업/규모별 비교, 다중응답 빈도); 인과모형 없음(저자 스스로 명시)\\

\citet{kim-2025-ai-adoption-firm-performance} &
통계청 기업활동조사(2017--2023), 상시근로자 50인이상$\cdot$자본금 3억이상 균형패널 약 7,000개 기업(약 42,000 관측치) &
TWFE(기업·연도·지역·산업FE, 기업단위 클러스터 표준오차)+사건연구(parallel trends 검증); TFP는 Levinsohn-Petrin(2003) 추정\\

\citet{lee-2025-ai-adoption-determinants-performance} &
사업체패널조사(WPS) 1--10차(결정요인 7--10차, N=6,095); 2023년(10차) OECD(2023) 동일문항 AI설문 최초 추가, 30인이상 사업체 &
결정요인은 CRE(Wooldridge 2005) 패널프로빗(TOE이론); 성과효과는 TWFE+Callaway and Sant'Anna(2021) 다시점DID(코호트별 ATT)\\

\citet{nam-2026-ai-macroeconomic-impact-analysis} &
국가데이터처 기업활동조사 패널(2017--2023, 연속 5개연도 이상); PSM은 2018--23 AI최초도입기업 대 매칭된 미도입기업(기업단위 부분만) &
leave-one-out AI도입률을 도구변수로 한 패널IV(Kleibergen-Paap rk F 검정)+PSM-DID(이벤트시간더미 k=0,1,2,3); 종속변수 ln(1인당매출)$\cdot$ln(상용근로자)$\cdot$ln(매출)\\

\citet{seo-2025-ai-adoption-characteristics-business-reports} &
금융감독원 DART 사업보고서 "사업의 내용" 전수 텍스트(2010--2024, 38,522건) &
AI키워드 1회이상 등장 여부로 도입기업 간접식별하는 텍스트마이닝(정규표현식 전처리); 도입률$\cdot$문서길이$\cdot$어휘다양성 기술통계, 인과모형 없음\\

\citet{spri-2024-ai-industry-survey} &
국가승인통계(제127016호), AI사업영위 상시근로자 1인이상 기업 2,517개 전수조사(기준일 2023.12.31) &
전수조사 기술통계(매출$\cdot$인력$\cdot$수출$\cdot$R\&D투자$\cdot$도입여부 항목별 집계, 전년대비 비교); 인과모형 없음\\

\citet{bick-2026-mind-the-gap-ai-adoption} &
자체 국제 근로자서베이(미국+유럽 6개국, Bilendi패널, 2025.5--6/2026.1--2 두 웨이브, N=55,415)+EU-ICT-Firm(32개국 157,000개 기업)+US BTOS+유럽 29개국 산업패널 &
Oaxaca-Blinder 분해로 미--유럽 도입률격차를 구성요인(인구$\cdot$산업$\cdot$직종$\cdot$규모)과 계수차이(고용주장려 등)로 분해; 국가$\times$산업 고정효과 패널회귀(인과관계 불주장 명시)\\

\citet{bick-2026-what-work-does-generative} &
Real-Time Population Survey(Qualtrics 온라인패널, 미국 18--64세) 2025.8$\cdot$11/2026.2$\cdot$5 4개 웨이브 풀링(웨이브당 약 5,000명), O*NET DWA(2,040개) 채택 신규모듈; 비교용 Anthropic/OpenAI/Microsoft 챗로그 &
직업$\cdot$과업단위 채택지수 구축; 노출점수(Eloundou $\alpha/\beta/\zeta$, Brynjolfsson, Webb, Felten)와 실제채택률 가중회귀; 학습채널은 과거경험 직업FE회귀+OAI 도구변수 2SLS+직장-비직장 파급 2SLS\\

\citet{cheon-2025-ai-exposure-vs-ai-adoption} &
한국 AI대표기업 51개(과기정통부$\cdot$NIPA)+AI스타트업 339개(한국인공지능협회) 중 355개 서비스텍스트; 한국고용직업분류 451개 직업 과업텍스트(6,133개 과업, 45,000개 DWA) &
Fenoaltea et al.(2024) AISE 방법론을 한국에 적용, GPT로 기업서비스-직업과업 텍스트매칭(예/아니오 판정 비율)해 AIFE 산출; AIOE$\cdot$AIOE\_by\_GPT와 피어슨 상관분석, 성별$\cdot$연령$\cdot$기업규모$\cdot$근속$\cdot$임금계층$\cdot$산업별 4분면 교차분석\\

\citet{kiet-2024-ai-labor-market-industry} &
통계청 기업활동조사 패널(2017--2021), 전체 15,993개 기업(AI도입 975개, 6.10\%) (기업단위 DID 부분만) &
표준DID(AI$\times$Year)+Callaway and Sant'Anna(2021) 다기간DID(코호트별 ATT); 종속변수 로그 종사자수/매출/손익; 강건성은 산업$\cdot$매출규모$\cdot$시차 재추정\\

\citet{chang-2024-ai-development-employment-effects} &
4장: 사람인$\cdot$잡코리아 구인공고(2017--2023)+고용보험DB(132,608개 사업체, AI도입 5,361개); 5장: 기업활동조사(2017--2022, 횡단면풀링 N=64,492/균형패널 N=44,157) (AI도입률 부분만) &
4장 OLS(로그 신규채용, AI도입더미+통제); 5장 Pooled OLS/IV(2SLS)/IV+Entropy Balancing(Czarnitzki et al.\ 2023)/패널FE$\cdot$RE$\cdot$HTE 병행, 도구변수는 AI인프라$\cdot$산업내 AI활용도\\

\citet{noh-2025-ai-labor-coexistence} &
3장: KBS$\cdot$한국노동연구원$\cdot$경총 공동설문(사업체N=110$\cdot$근로자N=3,227 병렬); 4장: 사무/관리직 임금근로자설문 N=500; 5--9장: 산업별 질적 인터뷰 (3--9장 설문 부분만) &
기업-근로자 병렬문항 비교(기술통계); 접근성$\cdot$과정$\cdot$결과 3차원 유형화; 질적 인터뷰 해석적 분석. DID/IV 등 인과모형 미사용(서술적$\cdot$탐색적 분석 명시)\\

\citet{noh-2025-ai-based-manufacturing-innovation-employment} &
고용보험DB(피보험자+사업장+구인공고, 2021.1--2024.12); 제조대기업 R\&D엔지니어설문 N=90+인터뷰 7개사(3장); 생산현장설문 N=102+인터뷰 7개사(4장) &
Babina et al.(2020) 방식(구인공고 AI숙련요구 시 도입기업 판정)으로 기업 AI도입 식별; 2022.11(ChatGPT출시)=100 지수 기준 준-사건연구식 추세비교; 3$\cdot$4장 기술통계$\cdot$교차분석 중심, 인과추론모형 미사용\\
\end{longtable}

% ============================================================
\section{실제 사용(Usage) 데이터 기반 실증연구}
\label{sec:usage}
% ============================================================

\subsection{데이터 출처 유형}

실사용 데이터 기반 연구의 출처는 채택 연구보다 훨씬 좁게 두 갈래—AI 기업의 대화 로그와 자기보고 서베이—로 수렴한다.

\begin{itemize}[leftmargin=*]
  \item \textbf{Anthropic Economic Index(AEI): Claude 대화 로그}: \citet{handa-2025-which-economic-tasks-are}가 처음으로 Claude.ai 대화 100만 건(2024.12--2025.1)을 프라이버시 보존 분석시스템(Clio)으로 O*NET 과업에 매핑한 이후, \citet{tamkin-2025-estimating-ai-productivity-gains-from}, \citet{massenkoff-2026-labor-market-impacts-of-ai}, \citet{massenkoff-2026-cadences}, \citet{massenkoff-2026-learning-curves}, \citet{appel-2025-uneven-geographic-and-enterprise-ai}, \citet{appel-2026-economic-primitives}, \citet{fan-2026-what-countries-use-ai-and}, \citet{fan-2026-aggregate-gains-from-ai}가 2025.8부터 2026.6까지 거의 매달 새로운 스냅샷(각 100만 건 내외)을 발행하며 시계열을 축적하고 있다. 원자료(개별 대화 내용)는 Anthropic 외부 연구자가 직접 접근할 수 없고, Anthropic 자체 연구팀이 발행하는 \emph{가공된 집계 보고서}만이 공개된다는 공통된 자료적 특성을 갖는다.
  \item \textbf{자기보고 가계조사(반복 단면)}: 한국은행 조사국이 수행하는 \citet{bok-2025-rapid-ai-diffusion-productivity-effects}(2025.5--6)와 \citet{bok-2026-ai-adoption-productivity-effects}(동일 표본설계 계승)는 한국리서치 웹패널을 통해 지역별고용조사 기준으로 층화된 대표표본(N=5,512)에게 업무 목적 AI 활용 경험·시간 절감률·업무처리량 변화를 자기보고로 직접 질문한다. Claude 대화 로그와 달리 특정 AI 서비스 이용자로 표본이 국한되지 않고 국가 대표표본이라는 장점이 있으나, 자기보고 특유의 편향(사회적으로 바람직한 응답, 절감률 과대추정 가능성)을 안고 있다.
  \item \textbf{실험실 실험}: \citet{min-2025-productivity-gains-generative-ai}는 온라인 실험 참가자(N=553)를 대상으로 생성형 AI 활용 여부를 직접 처치변수로 조작하는 실험을 수행한다. 관측 자료가 아니라 연구자가 직접 생성한 자료라는 점에서 유일한 사례다.
\end{itemize}

\citet{bick-2026-what-work-does-generative}의 서베이 기반 채택지수(\S\ref{sec:adoption}에서 다룸)는 챗로그 데이터와의 상관을 직접 검증한다는 점에서 채택과 실사용 두 갈래를 잇는 가교 역할을 한다.

\subsection{방법론 유형}

실사용 데이터의 방법론은 자료 자체의 짧은 역사(2023년 이후) 때문에 채택 연구보다 인과추론이 약하고, 대신 지표 구축과 대규모 기술통계(descriptive at scale)에 방법론적 기여가 집중된다.

\begin{itemize}[leftmargin=*]
  \item \textbf{과업 매핑 분류기(Clio)}: \citet{handa-2025-which-economic-tasks-are}는 대화를 O\*NET 과업문에 계층적으로 매핑하고 협업패턴 5유형(자동화 vs.\ 증강)으로 분류하는 프라이버시 보존 분류기를 방법론적 기여로 제시한다.
  \item \textbf{자기일관성 기반 시간절감 추정 + Hulten 집계}: \citet{tamkin-2025-estimating-ai-productivity-gains-from}는 Claude 자신에게 "AI 유무 시 소요시간"을 추정시키고(자기일관성 로그상관 $r=0.89$--$0.93$, JIRA 외부벤치마크 $\rho=0.44$--$0.50$으로 검증), 시간절감률 $=1-\text{time}_{\text{with}}/\text{time}_{\text{without}}$을 Hulten 정리로 경제 전체에 집계한다. \citet{appel-2026-economic-primitives}는 이를 과업 성공률로 할인해 재추정한다.
  \item \textbf{이론노출 게이팅 + 준실험 DID}: \citet{massenkoff-2026-labor-market-impacts-of-ai}는 Eloundou 이론노출을 실제 Claude 사용량 임계치로 게이팅한 "관측노출(observed exposure)" 지표를 구축하고, 이를 CPS 패널과 결합한 이중차분(ChatGPT 출시 전후)으로 청년 고용에 대한 인과효과를 검증한다—실사용 연구 중 인과추론 강도가 가장 높은 사례다.
  \item \textbf{학습곡선·산출물 분류 회귀}: \citet{massenkoff-2026-learning-curves}는 이용 경력(tenure)별 과업성공률 차이를 직업·모델·용도·국가 고정효과를 단계적으로 통제한 회귀로, \citet{massenkoff-2026-cadences}는 산출물 30여개 범주 신규분류기와 토큰-임금 로그-로그 회귀, 서베이-실사용데이터 프라이버시 보존 연계로 각각 분석한다.
  \item \textbf{지리적 분해 지수}: \citet{appel-2025-uneven-geographic-and-enterprise-ai}는 지역 사용점유율을 근로연령인구 점유율로 나눈 AI Usage Index(AUI)를 신규 정의해 국가·주 간 격차를 지수화한다.
  \item \textbf{중력모형(PPML)·집중도지수}: \citet{fan-2026-what-countries-use-ai-and}는 사용강도(PPML)와 확산폭(로그HHI, OLS)을 국가 특성에 회귀하고, \citet{fan-2026-aggregate-gains-from-ai}는 보건경제학의 concentration index를 응용한 AI집중지수(ACI)로 임금분포 내 이득 편중도를 측정한다.
  \item \textbf{생산함수 기반 상한 추정 + 자기보고 비교(WLS)}: \citet{bok-2025-rapid-ai-diffusion-productivity-effects}는 콥-더글러스 생산함수로 잠재적 생산성효과의 상한을 계산하고, \citet{bok-2026-ai-adoption-productivity-effects}는 이를 실제 업무처리량 증가(자기보고)와 대비하는 WLS 회귀($Y_i=\alpha+\beta X_i+Z_i'\gamma+(X_i\times Z_i)'\delta+F_i'\theta+\varepsilon_i$)로 "생산성 단절"을 검증한다.
  \item \textbf{실험설계}: \citet{min-2025-productivity-gains-generative-ai}는 2$\times$2 요인설계와 Welch's t-test·Kolmogorov-Smirnov·Wilcoxon Rank-Sum 검정을 병행한다.
\end{itemize}

\subsection{연구별 요약}

\begin{longtable}{L{2.5cm}L{6.4cm}L{6.4cm}}
\caption{실제 사용(Usage) 기반 실증연구: 데이터 출처와 방법론}\label{tab:usage}\\
\toprule
\textbf{문헌} & \textbf{데이터 출처} & \textbf{방법론}\\
\midrule
\endfirsthead
\multicolumn{3}{l}{\small (표 \thetable\ 계속)}\\
\toprule
\textbf{문헌} & \textbf{데이터 출처} & \textbf{방법론}\\
\midrule
\endhead
\bottomrule
\endfoot

\citet{handa-2025-which-economic-tasks-are} &
Claude.ai Free/Pro 대화(엔터프라이즈$\cdot$API 제외), 2024.12--2025.1 수집(7일 스냅샷), 과업수준 분석표본 100만건$\cdot$전체분석 400만건 이상(최소 5개 계정$\cdot$15건 대화 프라이버시 기준) &
Clio(프라이버시보존 분석시스템)로 대화를 O*NET 약 2만개 과업문에 계층매핑; 5가지 협업패턴(Directive/Feedback Loop=자동화, Task Iteration/Learning/Validation=증강) 분류(인간검증 병행); 임금$\cdot$Job Zone별 사용률$\cdot$모델간 비교\\

\citet{tamkin-2025-estimating-ai-productivity-gains-from} &
Claude.ai(Free/Pro/Max) 대화전사록 10만건(프라이버시보존 추출); 검증용 오픈소스저장소 JIRA티켓 1,000건 &
Claude 자신에게 AI유무 소요시간을 추정시켜(자기일관성 로그상관 $r$=0.89--0.93, JIRA벤치마크 $\rho$=0.44--0.50) 시간절감률=$1-\text{time}_{\text{with}}/\text{time}_{\text{without}}$ 산출; OEWS임금 매칭+Hulten정리로 경제전체 노동생산성/TFP 집계, 가속-병목과업 구분\\

\citet{massenkoff-2026-labor-market-impacts-of-ai} &
O*NET(약 800개 직업), Anthropic Economic Index Claude사용데이터(2025.8$\cdot$11), Eloundou et al.(2023) 이론노출등급 $\beta$; CPS패널(2016--), DOL 실업보험(UI) ETA203 &
"관측노출(Observed Exposure)" 신규지수($\tilde r_t=\mathbb{1}\{\text{WorkUsage}_t\ge100\}\times\mathbb{1}\{\beta_t\ge0.5\}\times\alpha_t$, 시간비중가중); AI노출-BLS고용전망 가중회귀; 노출상위25\% vs.\ 0\% 그룹 CPS DID(ChatGPT출시 기준), 청년(22--25세) 신규취업률 DID\\

\citet{massenkoff-2026-learning-curves} &
Claude.ai$\cdot$1P API 각 100만건(2026.2, 2025.11과 비교), BLS OEWS 2024.5 &
상위 10대 과업점유율$\cdot$과업가치(임금매핑) 시점간 비교; 직업별 Opus사용비중-시급 이분산회귀; 고tenure(6개월이상)/저tenure 3단계 회귀(단순$\to$과업FE$\to$모델$\cdot$용도$\cdot$국가FE 전체통제)로 과업성공률$\cdot$학습효과 식별\\

\citet{massenkoff-2026-cadences} &
Claude.ai/Cowork+1P API 대화(2026.4.10--6.10, 월/시간단위 집계); Anthropic Economic Index Survey(응답 약 9,700명, 설문-실사용데이터 프라이버시보존 연계) &
산출물(artifact) 30여개 범주 신규분류기; 토큰-임금 로그-로그 회귀(기하평균집계); AI자율성(1--5점) Claude.ai/Cowork vs.\ Claude Code 비교; 보고노출$\cdot$기대노출과 자동화사용비중의 관계 분석\\

\citet{appel-2025-uneven-geographic-and-enterprise-ai} &
Claude.ai Free/Pro 대화 100만건(2025.8.4--11); 1P API 고객트랜스크립트 100만건(2025.8); Census BTOS(인용) &
지리분해(150개국+미국 50개주+DC)로 Anthropic AI Usage Index(AUI=사용점유율/근로연령인구점유율) 신규정의; GDP-AUI 로그-로그회귀($R^2$=0.709); API vs.\ Claude.ai 협업모드$\cdot$비용-사용 탄력성(OLS, 편회귀 $-0.29$) 비교\\

\citet{appel-2026-economic-primitives} &
Claude.ai 소비자대화 100만건+1P API 100만건(2025.11.13--20, Opus4.5 직전); 인간연구자-분류기 비교+합성데이터로 검증 &
9개 신규분류기(과업복잡성$\cdot$인간/AI숙련$\cdot$용도$\cdot$AI자율성$\cdot$과업성공)로 5개 "경제원시지표" 구축; 실효AI커버리지=성공률$\times$시간비중 가중합; 과업임베딩 릿지회귀로 교육연수 예측; Hulten정리/CES집계를 성공률로 할인해 생산성 재추정\\

\citet{fan-2026-what-countries-use-ai-and} &
Anthropic Economic Index(2026.3.24 릴리스) Claude소비자대화 805,904건(2026.2.5--12, 100개국이상, IP기반 국가식별); WDI$\cdot$IMF AIPI 등 국가공변량 &
강도(인구백만명당대화수)는 PPML(구조적 영 포함); 확산폭(25개요청군집 로그HHI)은 OLS+HC1강건표준오차; Hidalgo et al.(2007) product space를 응용한 AI-diffusion space(RCA기반 근접도 네트워크)\\

\citet{fan-2026-aggregate-gains-from-ai} &
Anthropic Economic Index 5개웨이브(2025.1--2026.2, 웨이브당 약 100만대화); ILO 임금$\cdot$고용자료(ISCO-08 1자리, 86--117개국) &
노동비용등가(LCE)=과업별 절감시간$\times$임금$\times$사용비중 집계; AI집중지수(ACI, 보건경제학 concentration index 응용, $-1\sim+1$)로 임금분포 내 AI이득 편중도 측정; 웨이브FE+국가클러스터SE OLS\\

\citet{bok-2025-rapid-ai-diffusion-productivity-effects} &
한국은행 조사국 자체 가계조사, 통계청 지역별고용조사 기준 직업중분류$\cdot$연령$\cdot$성별 층화 다단계할당표집+사후층화가중, 한국리서치 웹패널 CAWI(2025.5.19--6.17, N=5,512) &
콥-더글러스 생산함수 기반 잠재적 생산성효과 $\approx(1-\alpha)\sum s_iw_i/\sum l_iw_i$(노동소득분배율 0.57, Acemoglu 2024 인용); 로봇협업확률은 현재협업자 특성기반 추정 후 비협업자 시뮬레이션 적용; WTP는 가상시나리오 설문\\

\citet{bok-2026-ai-adoption-productivity-effects} &
한국은행 조사국 가계조사(2025.5--6, 2025년 표본설계 계승, N=5,512, 회귀분석 N=1,801) &
"잠재적 생산성"(시간절감 기반 상한)과 "실현된 생산성"(업무처리량 자기보고 증가)을 구분; WLS(강건표준오차) $Y_i=\alpha+\beta X_i+Z_i'\gamma+(X_i\times Z_i)'\delta+F_i'\theta+\varepsilon_i$(교차항모형 조정 $R^2$=0.193); Bick et al.(2026a) 방법 원용한 산업별 AI활용률-생산성 산점도\\

\citet{min-2025-productivity-gains-generative-ai} &
컨슈머인사이트 모집 만19--60세 성인 온라인실험 참가자 602명중 553명 분석(2$\times$2 실험설계) &
순차 2과업(요약$\to$구매설득이메일)$\times$분업유무$\times$생성형AI(ChatGPT)활용유무 2$\times$2 실험(4개집단); ChatGPT기반 의미채점; Welch's t-test$\cdot$KS검정$\cdot$Wilcoxon Rank-Sum 병행으로 5개 가설 검증\\
\end{longtable}

% ============================================================
\section{데이터 출처와 방법론 비교: 채택 vs.\ 실사용}
\label{sec:compare}
% ============================================================

\subsection{데이터 출처의 체계적 차이}

두 갈래는 데이터 출처의 성격에서 뚜렷이 갈린다. 채택 연구는 한국의 경우 통계청 기업활동조사·WPS처럼 \textbf{여러 해에 걸쳐 반복되는 정부 행정·패널조사}에 크게 의존하며, 이 덕분에 TWFE·사건연구·다시점DID·IV 같은 준실험적 식별전략까지 나아갈 수 있는 자료적 기반을 갖춘다(\S\ref{sec:adoption} 참조). 반면 실사용 연구는 자료 자체가 2023년 이후에만 존재하는 신생 로그자료(Claude 대화 등)이거나, 마찬가지로 최근에야 시작된 반복 단면 서베이(한국은행 가계조사, 2025년 시작)에 국한되어 시계열이 짧다—이 때문에 \citet{massenkoff-2026-labor-market-impacts-of-ai}처럼 CPS라는 기존 반복 패널에 로그 기반 지표를 "이식"해야만 준실험적 인과추론이 가능해진다.

자기보고 방식(채택 다수, 한국은행 가계조사)의 공통 한계는 사회적으로 바람직한 응답 편향과 "도입"·"활용"의 모호한 자기 정의(어떤 도구를 얼마나 써야 "도입"인지 응답자마다 기준이 다를 수 있음)이다. 로그 기반 실사용 자료의 공통 한계는 표본이 특정 AI 서비스(Claude 등) 이용자로 국한되어 대표성이 제한되고, 원자료 자체를 외부 연구자가 볼 수 없어 발행 기관이 가공한 집계 보고서에 의존해야 한다는 점이다. \citet{cheon-2025-ai-exposure-vs-ai-adoption}의 AIFE는 이 두 한계 모두에서 벗어난 제3의 길을 보여준다—공급측(AI 기업의 실제 서비스 텍스트)에서 도출되므로 수요측 자기보고 편향이 없고, 355개 기업 전수(비표본) 자료라는 점에서 로그자료의 표본 편향 문제와도 다른 성격을 갖는다.

\subsection{방법론의 체계적 차이}

\begin{table}[h]
\centering
\begin{tabular}{L{2.6cm}L{6.0cm}L{6.0cm}}
\toprule
\textbf{식별전략 유형} & \textbf{채택 연구 예} & \textbf{실사용 연구 예}\\
\midrule
순수 기술통계 &
\citet{kim-2021-ai-adoption-status-major-industries}, \citet{spri-2024-ai-industry-survey}, \citet{seo-2025-ai-adoption-characteristics-business-reports} &
\citet{fan-2026-what-countries-use-ai-and}의 확산폭 지수, \citet{massenkoff-2026-cadences}의 산출물 분류\\
\midrule
회귀/고정효과(TWFE) &
\citet{kiet-2024-ai-adoption-firms-policy}, \citet{kim-2025-ai-adoption-firm-performance} &
\citet{bok-2026-ai-adoption-productivity-effects}의 WLS 교차항모형\\
\midrule
준실험(사건연구$\cdot$다시점DID) &
\citet{lee-2025-ai-adoption-determinants-performance}, \citet{kiet-2024-ai-labor-market-industry} &
\citet{massenkoff-2026-labor-market-impacts-of-ai}의 CPS 이중차분\\
\midrule
도구변수/구조모형 &
\citet{humlum-2019-robot-adoption-and-labor-market}의 동태적 구조모형, \citet{nam-2026-ai-macroeconomic-impact-analysis}의 패널IV &
(해당 사례 없음 --- 실사용 연구 중 도구변수를 쓴 사례는 \citet{bick-2026-what-work-does-generative}의 2SLS뿐이며 이는 채택지수 연구임)\\
\midrule
자기일관성/자체평가 추정 &
(해당 없음) &
\citet{tamkin-2025-estimating-ai-productivity-gains-from}의 Claude 자기추정+외부벤치마크 교차검증\\
\midrule
실험설계 &
(해당 없음) &
\citet{min-2025-productivity-gains-generative-ai}의 2$\times$2 실험\\
\bottomrule
\end{tabular}
\caption{식별전략 스펙트럼: 채택 vs.\ 실사용 연구}\label{tab:method-spectrum}
\end{table}

표\,\ref{tab:method-spectrum}이 보여주듯, 준실험·구조모형 수준의 강한 인과추론은 채택 연구 쪽에 집중되어 있다. 실사용 연구의 방법론적 기여는 대신 \emph{새로운 지표를 어떻게 대규모 로그에서 신뢰성 있게 구축하는가}(Clio 분류기, 자기일관성 검증, 경제원시지표 분류기)에 있으며, 이는 자료가 아직 짧아 전통적 계량경제학적 식별이 어렵다는 사정과, Anthropic 등 발행 주체가 스스로 방법론 혁신의 최전선에 있다는 사정이 겹친 결과로 해석된다.

% ============================================================
\section{결론: 본 프로젝트에 대한 시사점}
\label{sec:conclusion}
% ============================================================

본 정리는 채택과 실사용 두 갈래가 데이터 출처의 반복측정 가능성에서 근본적으로 다른 처지에 있음을 보여준다. 채택 연구는 정부 행정·패널조사라는 성숙한 자료 인프라 위에서 이미 준실험적 인과추론 단계에 도달했지만, 이분법적 더미(혹은 성숙도가 낮은 실태조사 단면)로는 활용의 \emph{강도}를 구분하지 못한다는 한계가 여전하다. 실사용 연구는 활용 강도를 가장 직접적으로 포착하지만, 자료의 짧은 역사와 특정 서비스 이용자로의 표본 국한이라는 이중의 대표성 문제를 안고 있다.

\citet{cheon-2025-ai-exposure-vs-ai-adoption}의 AIFE 구축 방법론—AI 공급기업의 실제 서비스 텍스트를 직업 과업과 LLM으로 매칭하는 방식—은 두 한계 모두를 피해가는 제3의 데이터 출처를 제시한다는 점에서, 향후 \texttt{data/proc/merged/} 설계 시 한국 특화 "실현 지표"의 구축 템플릿으로 검토할 가치가 있다. 동시에 \citet{massenkoff-2026-labor-market-impacts-of-ai}의 관측노출(observed exposure) 게이팅 방법—이론적 노출과 실제 로그 사용량을 하나의 지표로 결합해 기존 반복 패널(CPS)에 이식하는 방식—은 채택 연구가 보유한 반복 패널의 인과추론 역량과 실사용 로그의 세밀함을 동시에 살리는 벤치마크로 참고할 만하다. 향후 매칭 키(직업코드·산업코드·시점) 설계 시, 자기보고 기반 채택 지표(도입 정의의 모호성)와 로그 기반 실사용 지표(표본 대표성)가 서로 다른 방향의 편향을 가질 수 있음을 함께 점검해야 한다.

\clearpage
\printbibliography

\end{document}
